\documentclass[10pt,a4paper]{amsart}
\usepackage[margin=19mm]{geometry}
\usepackage{amsmath,amssymb,mathtools}
\usepackage[scr=rsfs,cal=euler]{mathalfa}
\usepackage{xfrac}
\usepackage[unicode,hidelinks,bookmarks=false]{hyperref}
\newcommand{\ie}{i.e.\ }
\newcommand{\cf}{cf.\ }
\newcommand{\resp}{respectively\ }
\newcommand{\Subsection}{\subsection}
\providecommand{\nobreakdash}{\nobreak\hbox{-}}
\setlength{\emergencystretch}{2em}
\multlinegap0pt
\usepackage{bm}
\usepackage{bbm}
\usepackage{dsfont}
\usepackage{xspace}
\usepackage{cancel}
\usepackage{mathtools}
\usepackage[shortlabels]{enumitem}
\usepackage{amsthm}
\makeatletter
\@ifundefined{theorem}{\newtheorem{theorem}{Theorem}}{}
\@ifundefined{lemma}{\newtheorem{lemma}[theorem]{Lemma}}{}
\@ifundefined{proposition}{\newtheorem{proposition}[theorem]{Proposition}}{}
\makeatother
\newcommand{\ee}{\quad\text{ and }\quad}
\newcommand{\steq}{\succ_{stab}}
\title[Characterizing $(d,h)$-elliptic stable irreducible curves]{Characterizing $(d,h)$-elliptic stable irreducible curves}
\author{Juliana Coelho, Renata Costa}
\date{Source: arXiv:2602.04865v1 (CC BY 4.0)}
\begin{document}
\maketitle

\noindent\emph{Source and licence.} Characterizing $(d,h)$-elliptic stable irreducible curves. Version arXiv:2602.04865v1 (CC BY 4.0), distributed under the Creative Commons Attribution 4.0 International licence, as recorded in the arXiv licence metadata for this exact version and shown on the abstract page https://arxiv.org/abs/2602.04865v1. Full rights notice: licenses/arxiv-2602.04865-CC-BY-4.0.txt.\par\smallskip

\noindent\textit{Editorial scope.} Counted unit: the Proposition whose source label is \texttt{prop:degreeofrestriction} in the article (source0.tex lines 499--503, source label \texttt{prop:degreeofrestriction}), delivered together with the article's own complete original proof of it (source0.tex lines 504--565, one \texttt{proof} environment of 62 lines). No mathematics is written, repaired or re-derived: statement and proof are the article's own text line for line. Required context retained uncounted: lem:3lemas (lines 463--473). Every definition, display and label of those lines is retained unchanged. Omitted from this package and not redistributed: 1--462, 474--498, 566--768. Nothing inside a retained excerpt is omitted or altered: every retained body is delivered exactly as the article's lines are written, so no omission receipt of any kind accompanies this package. Citations: the retained excerpts carry no citation command at all, so no bibliography entry of the article is transcribed and no reference is redistributed. Reference adaptation: none was needed and none was made; every reference inside the delivered unit resolves inside it, so no unresolved reference or citation is produced. Printed slips kept literal: no printed word, symbol, hypothesis or number of the counted statement or of its proof was altered. Layout and class: the article's own class is \texttt{article} with packages including amsmath, amsthm, amsfonts, amssymb, amscd, xy, graphicx, xcolor, enumerate, comment, hyperref, biblatex; the wrapper is the lane's pinned \texttt{amsart} editorial wrapper at 10pt on A4, which restates the theorem-like environments the retained text needs and adds the article's own macro definitions that the retained text needs (\texttt{\textbackslash ee}, \texttt{\textbackslash steq}), with extra wrapper packages (bm, bbm, dsfont, xspace, cancel, mathtools, enumitem). The delivered numbering is the unit's own. Assets: no retained span contains a figure environment, a graphics-inclusion command, a PDF-page inclusion, a table or a TikZ picture; no image file is copied into this package, the package declares no asset dependency and no image file is redistributed with it. Licence: version arXiv:2602.04865v1 of this article is distributed under the Creative Commons Attribution 4.0 International licence, as recorded in the arXiv licence metadata for this exact version and shown on the abstract page https://arxiv.org/abs/2602.04865v1. A full rights notice is delivered at licenses/arxiv-2602.04865-CC-BY-4.0.txt. Characters: every retained excerpt is pure ASCII, so the delivered engine, XeLaTeX with its default Latin Modern text font, reports no missing character.
\begin{lemma}\label{lem:3lemas}
Let $C$ be a nodal curve and let ${C}'$ be a nodal curve stably equivalent to $C$ with contraction map $\tau: {C}'\rightarrow C$. 
\begin{enumerate}[(a)]
\item If $Z_1$ and $Z_2$ are two connected subcurves of ${C}'$ such that $\tau(Z_1)= \tau(Z_2)$ is a subcurve of $C$, then $\tau(Z_1)= \tau(Z_2)= \tau(Z_1\cap Z_2)$.
\item If $C$ is irreducible then no two subcurves of ${C}'$ of positive genus have finite intersection.
\item Let $\pi: {C}' \rightarrow B$ be a $d$-sheeted  pseudo-admissible cover.
Assume that $C$ is irreducible and let $C_n$ be the normalization of $C$  and let $B_0 = \pi(C_n)$. Then every subcurve $Z$ of $B$ of positive genus must contain $B_0$.

In particular, if there exists a component $W$ of $B$ such that $g(W)\geq 1$, then $W=B_0$.
\end{enumerate}
\end{lemma}

\begin{proposition} \label{prop:degreeofrestriction}% \label{p3}
Let $C$ be an irreducible nodal curve and $C_n$ be its normalization. 
Let  $\pi:{C}'\rightarrow B$ be a $d$-sheeted pseudo-admissible cover, where ${C}'\steq C$. % is stably equivalent to $C$. 
If $g(B)\geq 1$, then $\pi|_{C_n}$ has degree $d$.
\end{proposition}

\begin{proof}
Set $B_0: = \pi(C_n)$ and 
let $e$ be the degree of the restriction  $\pi|_{C_n}$.  %$\pi_n := \pi|_{C_n} : C_n \rightarrow B_0$.
Assume $1\leq e \leq d-1$. 
Since $\pi$ is a map of degree $d$, then there exists a connected subcurve $Z$ of ${C}'$ not containing $C_n$ such that $C_n \cup Z\subset \pi^{-1}(B_0) $. 
If $e'$ is the degree of $\pi|_Z$, then $e'\geq 1$ and $e+e'\leq d$. 
We have two cases to consider.


If $g(B_0)\geq 1$ then $C_n$ and $Z$ must have positive genus as well, since they map to $B_0$. 
By Lemma \ref{lem:3lemas}, we then have that $C_n \cap Z$  is not finite
and,  since $C_n$ is irreducible, we must have $C_n\subset Z$ contradicting the choice of $Z$.

%
%\begin{figure}[!h]
%\centering
%\includegraphics[width=0.5\linewidth]{Modelo TD PGMAT-UFF/Mapa_.png}
%\caption{Case 2}
%\label{fig:case2}
%\end{figure}

Now assume $g(B_0)=0$. Then, by Lemma \ref{lem:3lemas},
every irreducible component of $B$ has genus 0  and, since $g(B)\geq 1$, there exists a rational cycle $E=(E_1,\ldots,E_s)$ in $B$ which  must contain $B_0$, say $E_1=B_0$. 
%Let $Z$ be the connected component of $\pi^{-1}(B\setminus B_0)$ containing the rational cycle $E$. 
Let
$q_1,p_1, \ldots, p_{s-1},q_2\in B$ be such that
$$E_1 \cap B_0 = {q_1},
\quad E_s \cap B_0 = {q_2},
\quad \text{ and } \quad 
E_{i}\cap E_{i+1}= {p_{i}},
\quad \text{for } i =1, \ldots, s-1. $$

Let $q^1_1 \in C_n$ and $q_1^2 \in Z$ such that
$$\pi(q_1^1)= \pi(q_1^2)= q_1$$
Since $q_1$ is a node and $\pi$ is pseudo-admissible, then $q_1^1$ and $q_1^2$ are also nodes. Therefore, there exist rational components $E_1^1$ and $E_1^2$ of $C'$  such that
$$q_1^1 = C_n \cap E_1^1\ee q_1^2= Z \cap E_1^2.$$


Note that we must have $\pi(E_1^1)= \pi(E_1^2)= E_1$. 
Furthermore, since $E_1$ has another node, namely $p_1$, then $E_1^1$ and $E_1^2$ also have nodes, say $p_1^1$ and $p_1^2$, respectively, such that
$$\pi(p_1^1)= \pi(p_1^2)= p_1. $$
Consequently, there exist rational components $E_2^1$ and $E_2^2$  of $C'$ such that
$$p_1^1 = E_1^1 \cap E_2^1\ee p_1^2= E_1^2 \cap E_2^2$$
and again
we have, $\pi(E_2^1)=\pi(E_2^2)=E_2$.


Prceeding in this manner, for $j=1,2$ we obtain rational components 
%$E_1^1,E_2^1,\ldots, E_{s-1}^1$ and
$E_1^j,E_2^j,\ldots, E_{s}^j$  of $C'$
such that  $\pi(E_i^j)=E_i$
and
$E^j_i$ meet $E^j_{i+1}$ at $p^j_{i+1}$ for $i=1,\ldots,s-1$.
Lastly, since $E_{s}$ meets $B_0$ at the node $q_2$, then 
$ E_{s}^1$ meets $C_n$ and 
$ E_{s}^2$ meets $Z$.
Hence 
$$C_n \cup E_1^1 \cup \ldots \cup E_s^1 
\ee
Z\cup E_1^2\cup \ldots \cup E_s^2$$
are two subcurves of ${C}'$ of positive genus and having no common component, contradicting Lemma \ref{lem:3lemas}, thus showing that $e=d$.
 \end{proof}   

\end{document}
